\section{Limitations and Future Work}

Despite the high quality of OpenWebMath, we note several limitations and avenues for future works. First, due to the high cost of extracting data from all shards on Common Crawl, we were only able to run our pipeline once. Therefore, many of our choices are without empirical justification and we provide no ablation study. We also note that the nature of this particular type of dataset means that there are many subjective choices to be made. For instance, what counts as a mathematical document? What is a high-quality document? How do we choose the threshold for near-deduplication? For each of these, we chose several values and manually inspected a few examples to choose. Due to the cost constraints, there are also practical challenges with balancing cost with accuracy when filtering and extracting text. For instance, our prefilter reduces the number of HTML documents processed to under 1\% of the documents in Common Crawl, which may be too aggressive. We also note that OpenWebMath is an English-only dataset, which limits its applications for researchers and users who speak other languages. Finally, we note that OpenWebMath only contains the \textit{text} from math on the web, not associated figures, which can be important for solving mathematical problems \citep{openai2023gpt4}. Future work should focus on finding empirical answers to the questions of what constitutes good data, creating new, efficient filtering methodologies, and extracting images inline with math text.

\section{Text Extraction}
\label{appendix:text_extraction}

\begin{table}[t]
\begin{center}
\begin{tabular}{lrl}
\multicolumn{1}{c}{\bf Method}  &\multicolumn{1}{c}{\bf Runtime (s)} &\multicolumn{1}{c}{\bf Source Code Link} \\
\midrule
Resiliparse & 3.99 & \href{https://github.com/chatnoir-eu/chatnoir-resiliparse}{https://github.com/chatnoir-eu/chatnoir-resiliparse} \\
HTML-Text & 10.75 & \href{https://github.com/TeamHG-Memex/html-text}{https://github.com/TeamHG-Memex/html-text} \\
Inscripts & 19.14 & \href{https://github.com/weblyzard/inscriptis}{https://github.com/weblyzard/inscriptis} \\
BoilerPy & 24.94 & \href{https://github.com/jmriebold/BoilerPy3}{https://github.com/jmriebold/BoilerPy3} \\
jusText & 31.17 & \href{https://github.com/miso-belica/jusText}{https://github.com/miso-belica/jusText} \\
HTML2Text & 37.17 & \href{https://github.com/Alir3z4/html2text/}{https://github.com/Alir3z4/html2text/} \\
BeautifulSoup & 38.42 & \href{https://code.launchpad.net/beautifulsoup}{https://code.launchpad.net/beautifulsoup} \\
Trafilatura & 63.90 & \href{https://github.com/adbar/trafilatura}{https://github.com/adbar/trafilatura} \\
ExtractNet & 299.67 & \href{https://github.com/currentslab/extractnet}{https://github.com/currentslab/extractnet} \\
\end{tabular}
\end{center}
\caption{We measured the performance of various HTML text extraction tools on a dataset of 1k documents. Resiliparse was by far the most efficient, leading us to choose it for use in our pipeline.}
\label{table:extraction-methods-comparison}
\end{table}


\paragraph{Choice of Base Text Extractor} When considering which HTML text-extraction library to use, we considered the efficiency, customization, and existing boilerplate reduction methods for each option. The most commonly used option, using WET files extracted by Common Crawl, was not an option since they do not deal with \LaTeX{} correctly and offer no customization. Other options such as jusText \citep{justext}, used in The Pile \cite{gao2020pile}, removed boilerplate too aggressively, leading to sections containing math to be discarded. Likewise, Trafilatura \citep{barbaresi-2021-trafilatura}, which was used in RefinedWeb \citep{refined-web}, had poor efficiency. We decided to go with Resiliparse \citep{bevendorff:2018} due to its balanced boilerplate removal, fast runtime, and efficient Common Crawl parsing tools. \autoref{table:extraction-methods-comparison} shows the full results for our comparison.

\begin{table}[t]
\begin{center}
\begin{tabular}{lr}
\multicolumn{1}{c}{\bf Math Format} & \multicolumn{1}{c}{\bf Percentage of Documents} \\
\midrule
Found at least one instance of math  & 91.42\% \\
MathJax with delimiters (inline) & 50.27\% \\
MathJax with delimiters (display) & 23.37\% \\
Math found in images & 6.96\% \\
\texttt{.math-container} & 3.94\% \\
MathML code & 3.28\% \\
\texttt{<annotation>} withing \texttt{<math>} tags & 2.35\% \\
\texttt{<mathjax>} tags & 2.24\% \\
\texttt{align} environments & 1.72\% \\
\texttt{equation} environments & 1.18\% \\
within \texttt{<script>} tags & 1.01\% \\
\texttt{alttext} property of \texttt{<math>} tags & 0.24\% \\
\end{tabular}
\end{center}
\caption{Frequencies of different types of \LaTeX{} found in OpenWebMath. The most common format of \LaTeX{} found in Common Crawl is MathJax, which uses user-defined delimiters to denote math equations. Second most common is \LaTeX{} code within either the URL or \texttt{alt} text of an \texttt{img} tag.}
\label{table:latex-types}
\end{table}

\paragraph{\LaTeX{} Extraction} \LaTeX{} code comes in many forms throughout Common Crawl HTML files. We employed an iterative process to refine our extraction rules. First, we filtered shards of Common Crawl for documents that contain the string \texttt{\textbackslash frac}. Then, we filtered those documents to find those which our extraction code found no extractable \LaTeX{}. Then, we refined our code to include additional sources of math until we were confident that we had reasonable support for all formats of \LaTeX{} in HTML documents. \autoref{table:latex-types} shows the breakdown of different common types of \LaTeX{} found in HTML documents. 

We note that most of the \LaTeX{} in OpenWebMath and across the internet is encoded using MathJax, which presents a challenge. The majority of MathJax documents use dollar sign delimiters, but most dollar signs on the web do not delimit \LaTeX{} equations. This leaves us with a few options:

\begin{itemize}
    \item Detect the use of the MathJax script in the HTML file. If the script is imported, treat dollar signs as \LaTeX{} code.
    \item Detect common \LaTeX{} commands in between dollar signs. If they are present, treat dollar signs as \LaTeX{} code.
    \item Use the \textit{MathScore} classifier to determine whether the page looks like it is talking about math. If so, treat dollar signs as \LaTeX{} code.
\end{itemize}

The first option is not always accurate since the MathJax javascript code may be nested inside of another import or named differently depending on the website. The latter two options make up for many of these cases, but can fail to detect edge cases where math equations are present but the surrounding text does not indicate that the document is mathematical. We suspect Minerva \citep{lewkowycz2022solving} gets around this issue by using HTML documents where javascript code has already been executed, in which case MathJax is converted from delimited text to explicit HTML tags that are easy to detect.

\begin{table}[t!]
    \centering
    \begin{tabular}{ccccccc}
        \toprule
        Model Size & Layers & Model Dim & Heads & Learning Rate & Batch Size \\
        \midrule
        1.4 B & 24 & 2048 & 16 & $2.0 \times 10^{-4}$ & 1M \\
        \bottomrule
    \end{tabular}
    \caption{Model Hyperparameters. We use the same architecture and hyperparameters, other than batch size, as Pythia 1.4B \citep{pythia}.}
    \label{table:model_hyperparameters}
\end{table}

\begin{table}[t!]
\caption{List of Math Keywords used in the prefiltering stage.}
\label{table:math-keywords}
\begin{center}
\begin{tabular}{l}
\multicolumn{1}{c}{\bf Math Keywords} \\
\midrule
\texttt{MathJax} \\
\texttt{mathjax} \\
\texttt{<math} \\
\texttt{math-container} \\
\texttt{katex.min.css} \\
\texttt{latex.php} \\
\texttt{codecogs} \\
\texttt{tex.cgi} \\
\texttt{class="tex"} \\
\texttt{class='tex'} \\
\end{tabular}
\end{center}
\end{table}

\section{Interplay Between Extraction and Filtering}

In prior works, we noticed many cases where suboptimal HTML text extractors were used and yet text quality remains high in the dataset. This is due to the interplay between extraction and filtering. Specifically, if a text extractor fails to extract the main text, gets the formatting wrong, or includes too much boilerplate in the extraction, then both the classification and perplexity filters can filter out such examples. This can lead to subtle biases in the dataset, where specific poorly-extracted websites are excluded entirely even though they do contain high quality content. In the case of making a mathematical dataset, failure to extract and deal with inline \LaTeX{} code properly can hurt perplexity scores and lead to these documents being filtered out. We suggest practitioners tune their text extraction pipeline on a diverse set of documents before applying filtering to avoid this bias.

\section{Model Hyperparameters}

We trained models on 14.7B tokens using the LLaMA \citep{llama2} tokenizer and the architecture described in Pythia \citep{pythia}. We train the model using the GPT-NeoX library \citep{gpt-neox-library} on 8 A100 80GB GPUs. Exact hyperparameters can be found in \autoref{table:model_hyperparameters}.

\clearpage
\section{Datasheet}

We provide a datasheet for OpenWebMath, following the framework in \citet{gebru2021datasheets}.
\begin{longtable}{p{6cm}|p{6cm}}
    \toprule
    \multicolumn{2}{c}{\textsc{\textbf{Motivation}}} \\
    \midrule
    \textbf{For what purpose was the dataset created?} & The dataset was created to enable the training of large language models on mathematical texts, in order to improve their mathematical reasoning capabilities.\\ \midrule
    \textbf{Who created the dataset and on behalf of which entity?} & The dataset was created by the authors of this work. \\ \midrule
    \textbf{Who funded the creation of the dataset?} & Resources used in preparing this research were provided, in part, by the Province of Ontario, the Government of Canada through CIFAR, Fujitsu Limited, and companies sponsoring the Vector Institute for Artificial Intelligence (\url{www.vectorinstitute.ai/partners}). Computing resources for model training were provided by EleutherAI and Brigham Young University. \\ \midrule
    \textbf{Any other comment?} &  None. \\ \midrule
    \multicolumn{2}{c}{\textsc{\textbf{Composition}}} \\ \midrule
    \textbf{What do the instances that comprise the dataset represent?} & The instances are text documents extracted from mathematics-related webpages from Common Crawl. \\ \midrule
    \textbf{How many instances are there in total?} & In total, OpenWebMath contains 6.3 million documents. \\ \midrule
    \textbf{Does the dataset contain all possible instances or is it a sample (not necessarily random) of instances from a larger set?} & OpenWebMath doesn't contain all instances of text extracted from mathematics-related webpages from Common Crawl, as our filters can miss a non-zero proportion of such webpages. However, we expect OpenWebMath to contain most of them. \\ \midrule
    \textbf{What data does each instance consist of?} & Each instance consists of plain text and metadata including the source URL, the snapshot date, and other extraction parameters. \\ \midrule
    \textbf{Is there a label or target associated with each instance?} & No. \\ \midrule
    \textbf{Is any information missing from individual instances?} & No. \\ \midrule
    \textbf{Are relationships between individual instances made explicit?} & No. \\ \midrule
    \textbf{Are there recommended data splits?} & No. \\ \midrule
    \textbf{Are there any errors, sources of noise, or redundancies in the dataset?} & Yes, a small portion of the documents from OpenWebMath are not related to mathematics, or contain bad quality content. \\ \midrule
    \textbf{Is the dataset self-contained, or does it link to or otherwise rely on external resources?} & The dataset is entirely self-contained. \\ \midrule
    \textbf{Does the dataset contain data that might be considered confidential?} & No. \\ \midrule
    \textbf{Does the dataset contain data that, if viewed directly, might be offensive, insulting, threatening, or might otherwise cause anxiety?} & The data is filtered for quality and we do not expect that this content will be offensive, but since our filters may be imperfect we make no guarantees. \\ \midrule
    \multicolumn{2}{c}{\textsc{\textbf{Collection}}} \\ \midrule
    \textbf{How was the data associated with each instance acquired?} & The data was acquired by processing data from Common Crawl. \\ \midrule
    \textbf{What mechanisms or procedures were used to collect the data?} & We refer to the CommonCrawl website (\href{https://commoncrawl.org/}{commoncrawl.org}) for details on how they collect data. \\ \midrule
    \textbf{If the dataset is a sample from a larger set, what was the sampling strategy?} & We use all data from Common Crawl that was available before May 2023. \\ \midrule
    \textbf{Who was involved in the data collection process and how were they compensated?} & Keiran Paster and Marco Dos Santos collected the data and were compensated by their respective graduate programs. \\ \midrule
    \textbf{Over what timeframe was the data collected?} & OpenWebMath uses shards of CommonCrawl gathered between 2013 and 2023.   \\ \midrule
    \textbf{Were any ethical review processes conducted?} & No. \\ \midrule
    \multicolumn{2}{c}{\textsc{\textbf{Preprocessing}}} \\ \midrule
    \textbf{Was any preprocessing/cleaning/labeling of the data done?} & Yes. See \autoref{sec:filtering} for details. \\ \midrule
    \textbf{Was the “raw” data saved in addition to the preprocessed/cleaned/labeled data?} & Yes. \\ \midrule
    \textbf{ Is the software that was used to preprocess/clean/label the data available?} & Yes. See supplementary materials. \\ \midrule
    \multicolumn{2}{c}{\textsc{\textbf{Uses}}} \\ \midrule
    \textbf{Has the dataset been used for any tasks already?} & Yes, the data was used to train 1.4B parameter language models in \autoref{sec:models} \\ \midrule
    \textbf{Is there a repository that links to any or all papers or systems that use the dataset?} & No. \\ \midrule
    \textbf{What (other) tasks could the dataset be used for?} & We primarily envision that OpenWebMath could be useful for language model pretraining, finetuning, and evaluation. \\ \midrule
    \textbf{Is there anything about the composition of the dataset or the way it was collected and preprocessed/cleaned/labeled that might impact future uses?} & It is possible that the filtering stage of the project discarded valuable documents, such as those not written in English. This makes OpenWebMath suboptimal for creating mathematical models in other languages. \\ \midrule
    \textbf{Are there tasks for which the dataset should not be used?} & Any tasks which may considered irresponsible or harmful. \\ \midrule
    \multicolumn{2}{c}{\textsc{\textbf{Distribution}}} \\ \midrule
    \textbf{Will the dataset be distributed to third parties outside of the entity on behalf of which the dataset was created?} & Yes, the dataset will be available on the Hugging Face Hub for NLP practitioners. \\ \midrule
    \textbf{How will the dataset will be distributed?} & We will distribute the dataset on the Hugging Face Hub \\ \midrule
    \textbf{When will the dataset be distributed?} & The dataset will be available when the paper is made public. \\ \midrule
    \textbf{Will the dataset be distributed under a copyright or other intellectual property (IP) license, and/or under applicable terms of use (ToU)?} & The public extract is made available under an ODC-By 1.0 license; users should also abide to the CommonCrawl ToU: \href{https://commoncrawl.org/terms-of-use/}{https://commoncrawl.org/terms-of-use/}. \\ \midrule
    \textbf{Have any third parties imposed IP-based or other restrictions on the data associated with the instances?} & Not to our knowledge. \\ \midrule
    \textbf{Do any export controls or other regulatory restrictions apply to the dataset or to individual instances?} & Not to our knowledge. \\ \midrule
    \multicolumn{2}{c}{\textsc{\textbf{Maintenance}}} \\ \midrule
    \textbf{Who will be supporting/hosting/maintaining the dataset?} & The dataset will be hosted on the Hugging Face Hub. \\ \midrule
    \textbf{How can the owner/curator/manager of the dataset be contacted?} & \texttt{keirp@cs.toronto.edu} \\ \midrule
    \textbf{Is there an erratum?} & No. \\ \midrule
    \textbf{Will the dataset be updated?} & No. \\ \midrule
    \textbf{If others want to extend/augment/build on/contribute to the dataset, is there a mechanism for them to do so?} & No. \\ \bottomrule
    \caption{\textbf{Datasheet for OpenWebMath}, following the framework introduced by \citet{gebru2021datasheets}.}
    \label{tab:datasheet}
\end{longtable}
\clearpage
